\section{Projektphasen und Zeitplanung}

\subsection{Analyse und Planung}

\begin{itemize}[itemsep=0.2pt]
    \item Bestehende Datenstruktur in den GitHub-Repositories analysieren.
    \item Der Technologie-Stack musste besprochen und gewählt werden, sowie Einarbeitung in diesen.
    \item Die Anforderung an Datenpipline, API und Frontend mit dem Betrieb durchgehen und festlegen.
    \item Graphstruktur für die Aufgaben(Challenges) und Beziehungen(Abhängigkeiten) finalisieren da sie durch die Lern-Repositories teilweise gegeben sind.
\end{itemize}

\subsection{Datenpipline und Skripte}

\begin{itemize}[itemsep=0.2pt]
    \item Bash-Skripte zum Herunterladen, Aktualsieren oder zur erweiterung der Lern-Repositories.
    \item Kommentarblöcke und relevante Metadaten aus denn Repositories extrahieren.
    \item Erzeugung und Validierung der \texttt{graph-data.json} umsetzten.
    \item Implementierung der Daten in die Neo4j-Datenbank.
\end{itemize}

\subsection{Neo4j und API}

\begin{itemize}[itemsep=0.2pt]
    \item Das Neo4j-Datenmodell und Importverfahren per Skript umsetzten.
    \item Knoten(Challenges) und Beziehung(Abhängigkeiten) per Python im Skript strukturiert importieren.
    \item FastAPI und Neo4j-Client verbinden.
    \item API-Endpoints für Graph-, Listen-, Keywordsansicht der Challenges umsetzten.
    \item Implementierung von Admin-Endpoints, um per Zeitplan oder manuell die Challengedaten aktualisieren zu können.
    \item Validierung der Statusdaten um ihn im Frontend darstellen zu können.
\end{itemize}

\subsection{Frontend}

\begin{itemize}[itemsep=0.2pt]
    \item Moderne Frontend-Grundstruktur und Route-Navigation erstellen.
    \item API-Anbindung implememtieren.
    \item Graph-, Listen-, Keywordansicht der Challenges umsetzten.
    \item Hover-, Detailansicht und Filterung integrieren.
\end{itemize}

\subsection{Test und Qualiätssicherung}

\begin{itemize}[itemsep=0.2pt]
    \item Skripte, Status und JSON-Erzeugung testen.
    \item API-Endpoints und die Ausgabe testen.
    \item Datenbestand zwischen der JSON und Neo4jdatenbank vergleichen.
    \item Frontend-Anbindung und die Darstellungen prüfen.
    \item Erkannte Fehler korrigieren und funktionsfähig machen.
\end{itemize}

\subsection{Dokumentation und Übergabe}

\begin{itemize}[itemsep=0.2pt]
    \item Projektdokumenation und denn Antrag in LaTeX erstellt.
    \item Automatisierter Makefile und GitHub-Workflow Build
    \item Technische Übergabe und Besprechung mit dem Ausbilder durchführen
\end{itemize}